\section{量子物理现象}

\subsection{真空光场的模式与黑体辐射}

在边长 $l_1,l_2,l_3$ 的周期盒中，真空四势满足波动方程和 Lorenz 规范
\[
 \Box A^\mu=0,\qquad \partial_\mu A^\mu=0,
\]
允许波矢
\[
 k_i=\frac{2\pi N_i}{l_i},\qquad
 \omega=ck.
\]
对每个非零 $\vect k$ 只有两种相互正交的横向偏振具有物理电磁场。
沿 $k^\mu$ 的规范模式虽满足形式方程，却给出 $\vect E=\vect B=0$，不应误称为可观测“纵光子”。

把磁矢势 Fourier 展开为
\[
 \vect A(\vect r,t)=
 \sum_{\vect k,r}q_{\vect k,r}(t)\,
 \vect e_{\vect k,r}\e^{\ii\vect k\cdot\vect r},
\qquad r=1,2.
\]
利用 Parseval 公式，电磁场 Hamilton 量成为独立谐振子之和
\[
 H=\sum_{\vect k,r}
 \left(\frac{|p_{\vect k,r}|^2}{2m}
 +\frac{m\omega_{\vect k}^2|q_{\vect k,r}|^2}{2}\right).
\]
半径为 $k$ 的波矢球计数并乘两种偏振，单位体积角频率态密度为
\[
 \boxed{g(\omega)=\frac{\omega^2}{\pi^2c^3}}.
\]

经典能量均分给每个模式平均能量 $k_{\rm B}T$，得到 Rayleigh--Jeans 谱
\[
 u_{\rm RJ}(\omega)=\frac{k_{\rm B}T\omega^2}{\pi^2c^3},
\]
其高频积分发散。Planck 量子化令谐振子能级为
$E_n=n\hbar\omega$（真空零点项不影响热谱），配分函数与平均热能为
\[
 Z_\omega=\frac1{1-\e^{-\beta\hbar\omega}},\qquad
 \bar E_\omega=\frac{\hbar\omega}{\e^{\beta\hbar\omega}-1}.
\]
因此
\[
 \boxed{u(\omega)=
 \frac{\hbar\omega^3}{\pi^2c^3}
 \frac1{\e^{\hbar\omega/(k_{\rm B}T)}-1}}.
\]
低频极限回到 Rayleigh--Jeans 公式，高频呈指数衰减。

用
\[
 \int_0^\infty\frac{x^{s-1}}{\e^x-1}\,\dd x
 =\Gamma(s)\zeta(s)
\]
且 $\zeta(4)=\pi^4/90$，得到
\[
 u=a_{\rm R}T^4,\qquad
 a_{\rm R}=\frac{\pi^2k_{\rm B}^4}{15\hbar^3c^3},
\]
\[
 j=\frac{cu}{4}=\sigma_{\rm SB}T^4,\qquad
 \sigma_{\rm SB}=
 \frac{\pi^2k_{\rm B}^4}{60\hbar^3c^2}.
\]
Bernoulli 数的母函数
$t/(\e^t-1)=\sum B_nt^n/n!$
还给出低频渐近展开。

\begin{advtheorem}{$\pi\cot\pi z$ 的 Mittag--Leffler 展开}\label{advA6:thm:mittag-leffler}
对 $z\notin\mathbb Z$，
\[
 \boxed{\pi\cot\pi z=
 \frac1z+\sum_{n=1}^\infty
 \left(\frac1{z-n}+\frac1{z+n}\right)}.
\]
从而
\[
 \zeta(2m)=
 (-1)^{m+1}\frac{(2\pi)^{2m}}{2(2m)!}B_{2m}.
\]
\end{advtheorem}

\noindent\textit{证明。}
右端按 $\pm n$ 成对求和后局部一致收敛，并与 $\pi\cot\pi z$ 具有相同的一阶极点和留数。
令两者之差为 $h(z)$，则 $h$ 是整函数且周期为 $1$。
在矩形
$[-N-\tfrac12,N+\tfrac12]\times[-N,N]$ 上，
$\cot\pi z$ 在远离整数极点处有统一界；成对级数的尾项为 $O(N^{-1})$。
由最大模原理，令 $N\to\infty$ 得 $h$ 为有界整函数，故为常数。
取 $z\to\ii\infty$ 或比较 $z=1/2$ 可知常数为零。

再写
\[
 \pi z\cot\pi z=1-2\sum_{s=1}^\infty\zeta(2s)z^{2s},
\]
与 Bernoulli 展开
\[
 \pi z\cot\pi z=
 \sum_{s=0}^\infty
 (-1)^s\frac{(2\pi)^{2s}B_{2s}}{(2s)!}z^{2s}
\]
比较系数即得第二式。$m=1$ 给出 Basel 结果
$\zeta(2)=\pi^2/6$。

\begin{advexamplebox}{充满热辐射的弹性气球}\label{advA6:ex:radiation-balloon}
质量为 $m$ 的球形薄膜包住温度为 $T$ 的热辐射，外界为真空且无泄漏。
薄膜弹性等效于表面张力系数为 $\alpha$ 的肥皂泡。

\begin{enumerate}
\item 求平衡半径；
\item 若薄膜作各向同性快速径向小振动，辐射与薄膜来不及交换热量，求周期；
\item 判断这种热平衡在与热库接触时是否稳定。
\end{enumerate}
\end{advexamplebox}
\begin{advsolution}
光子气体能量密度与压强为
\[
 u=a_{\rm R}T^4,\qquad p=\frac u3.
\]
肥皂泡有内外两表面，Laplace 压强为 $4\alpha/R$，故
\[
 \frac{a_{\rm R}T^4}{3}=\frac{4\alpha}{R},
\qquad
 \boxed{R_0=\frac{12\alpha}{a_{\rm R}T^4}}.
\]

快速绝热运动中光子气体满足 $TR=\text{常量}$，所以
$U=a_{\rm R}T^4(4\pi R^3/3)=C/R$。
两表面的表面能为 $8\pi\alpha R^2$，
\[
 V_{\rm eff}(R)=\frac CR+8\pi\alpha R^2.
\]
平衡条件给 $C=16\pi\alpha R_0^3$，二阶导数
\[
 V_{\rm eff}''(R_0)=\frac{2C}{R_0^3}+16\pi\alpha
 =48\pi\alpha.
\]
薄膜所有质点径向速度均为 $\dot R$，动能为 $m\dot R^2/2$，故
\[
 \boxed{\omega_R=\sqrt{\frac{48\pi\alpha}{m}}},
\qquad
 \boxed{\mathcal T=2\pi\sqrt{\frac{m}{48\pi\alpha}}}.
\]

在平衡族上 $R_0\propto T^{-4}$，总能量
$U+8\pi\alpha R_0^2\propto T^{-8}$，因此
$C_{\rm eff}=\dd E/\dd T<0$。孤立体系可在力学上稳定振动，
但与无限热库作正反馈换热时热平衡不稳定。这区分了力学稳定性和热稳定性。
\end{advsolution}

\subsection{介质响应、因果性与散射}

Lorentz 振子模型
\[
 m\ddot x+m\Gamma\dot x+m\omega_0^2x=-eE
\]
给出
\[
 \chi(\omega)=
 \frac{\omega_p^2}{\omega_0^2-\omega^2-\ii\Gamma\omega},
\qquad
 \omega_p^2=\frac{n_eZe^2}{m\varepsilon_0}.
\]
写 $\chi=\chi_1+\ii\chi_2$，实部控制色散，虚部控制吸收。

\begin{advtheorem}{Kramers--Kronig 关系}\label{advA6:thm:kramers-kronig}
若线性响应满足因果性，$\chi(t<0)=0$，且高频足够快衰减，则
\[
 \boxed{\chi_1(\omega)=\frac1\pi
 \operatorname P\!\int_{-\infty}^{\infty}
 \frac{\chi_2(\omega')}{\omega'-\omega}\,\dd\omega'},
\]
\[
 \boxed{\chi_2(\omega)=-\frac1\pi
 \operatorname P\!\int_{-\infty}^{\infty}
 \frac{\chi_1(\omega')}{\omega'-\omega}\,\dd\omega'}.
\]
\end{advtheorem}

\noindent\textit{证明。}
因果卷积
\[
 P(t)=\varepsilon_0\int_0^\infty\chi(\tau)E(t-\tau)\,\dd\tau
\]
给出
\[
 \chi(z)=\int_0^\infty\chi(\tau)\e^{\ii z\tau}\,\dd\tau .
\]
当 $\operatorname{Im}z>0$ 时指数衰减，故 $\chi(z)$ 在上半平面解析。
对
$\chi(z')/(z'-z)$ 沿实轴和上半平面大半圆应用 Cauchy 公式，
高频衰减使半圆积分消失：
\[
 \chi(z)=\frac1{2\pi\ii}
 \int_{-\infty}^{\infty}\frac{\chi(\omega')}{\omega'-z}\,\dd\omega'.
\]
令 $z=\omega+\ii0^+$，使用 Sokhotski--Plemelj 公式
\[
 \frac1{x-\ii0^+}=\operatorname P\frac1x+\ii\pi\delta(x),
\]
再分别取实部、虚部，得到两式。

\begin{advexamplebox}{Lorentz 振子的因果脉冲响应}\label{advA6:ex:lorentz-laplace}
对
\[
 \chi(\omega)=
 \frac{\omega_p^2}{\omega_0^2-\omega^2-\ii\Gamma\omega}
\]
用 Laplace 变换求时域响应 $\chi(t)$，并检验因果性。
\end{advexamplebox}
\begin{advsolution}
单位脉冲驱动的响应满足
\[
 \ddot\chi+\Gamma\dot\chi+\omega_0^2\chi
 =\omega_p^2\delta(t),\qquad \chi(t<0)=0.
\]
单边 Laplace 变换给
\[
 \tilde\chi(s)=\frac{\omega_p^2}
 {s^2+\Gamma s+\omega_0^2}.
\]
欠阻尼 $\Gamma<2\omega_0$ 时，令
$\Omega=\sqrt{\omega_0^2-\Gamma^2/4}$，逆变换得
\[
 \boxed{\chi(t)=
 \Theta(t)\frac{\omega_p^2}{\Omega}
 \e^{-\Gamma t/2}\sin\Omega t}.
\]
它在 $t<0$ 严格为零，并在 $t=0^+$ 满足
$\dot\chi(0^+)=\omega_p^2$。令 $s=-\ii\omega+0^+$ 便恢复给定的
$\chi(\omega)$，其极点均在频率下半平面，因果性检查通过。
\end{advsolution}

加速电荷的 Larmor 辐射功率为
\[
 P_{\rm rad}=\frac{q^2a^2}{6\pi\varepsilon_0c^3}.
\]
Lorentz--Abraham 辐射反作用的非相对论形式
\[
 \vect F_{\rm rad}=m\tau\dot{\vect a},\qquad
 \tau=\frac{q^2}{6\pi\varepsilon_0mc^3}
\]
只是在变化时间尺度远大于 $\tau$ 时的有效近似。
把它当作精确三阶运动方程会出现 runaway 解或预加速；正确做法是作降阶近似
\[
 m\vect a\simeq\vect F_{\rm ext}+\tau\dot{\vect F}_{\rm ext},
\]
并只保留 $O(\tau)$。对电子 $\tau$ 约为 $6.3\times10^{-24}\,\mathrm s$。

谐振电子的辐射散射截面可写为
\[
 \sigma(\omega)=\frac{8\pi r_e^2}{3}
 \frac1{(\omega_0^2/\omega^2-1)^2+\omega^2\tau^2},
\qquad
 r_e=\frac{e^2}{4\pi\varepsilon_0mc^2}.
\]
高于束缚频率且仍远低于 $1/\tau$ 时趋于 Thomson 截面
$8\pi r_e^2/3$；低频极限 $\sigma\propto\omega^4$，产生 Rayleigh 散射；
在 $\omega\simeq\omega_0$ 附近出现共振增强。经典模型能说明形状，
但离散原子谱和跃迁强度需要量子理论。

\subsection{Stern--Gerlach 实验与自旋态}

非均匀磁场对磁矩的 $z$ 分量施力
\[
 F_z=\mu_z\frac{\partial B_z}{\partial z}.
\]
碱金属原子束分成两束，说明电子自旋投影只有
$S_z=\pm\hbar/2$。电子磁矩
\[
 \vect\mu_s=-g_e\frac{e}{2m}\vect S,\qquad
 g_e\simeq2.002319\ldots
\]
的负号来自电子负电荷。

\begin{center}\begin{minipage}{\linewidth}\centering
\centering
\begin{tikzpicture}[>=Latex,scale=.70]
 \begin{scope}[yshift=2.1cm]
  \node[draw,minimum width=1.3cm,minimum height=.65cm] (O1) at (-5,0){炉};
  \node[draw,minimum width=1.3cm,minimum height=.65cm] (Z1) at (-2.8,0){SG$_z$};
  \node[draw,minimum width=1.3cm,minimum height=.65cm] (Z2) at (1.1,.55){SG$_z$};
  \draw[->] (O1)--(Z1);
  \draw[->] (Z1.east)--(Z2.west) node[midway,above]{$z+$};
  \draw[dashed] (Z1.east)--(1.1,-.55) node[midway,below]{$z-$ 被挡};
  \draw[->] (Z2.east)--(4.7,.55) node[right]{$z+$};
  \draw[dashed] (Z2.east)--(4.7,-.15) node[right]{无 $z-$};
  \node[left] at (-6.1,0){(a)};
 \end{scope}
 \begin{scope}
  \node[draw,minimum width=1.3cm,minimum height=.65cm] (O2) at (-5,0){炉};
  \node[draw,minimum width=1.3cm,minimum height=.65cm] (Z3) at (-2.8,0){SG$_z$};
  \node[draw,minimum width=1.3cm,minimum height=.65cm] (X1) at (1.1,.55){SG$_x$};
  \draw[->] (O2)--(Z3);
  \draw[->] (Z3.east)--(X1.west) node[midway,above]{$z+$};
  \draw[dashed] (Z3.east)--(1.1,-.55) node[midway,below]{$z-$ 被挡};
  \draw[->] (X1.east)--(4.7,.95) node[right]{$x+$};
  \draw[->] (X1.east)--(4.7,.15) node[right]{$x-$};
  \node[left] at (-6.1,0){(b)};
 \end{scope}
 \begin{scope}[yshift=-2.15cm]
  \node[draw,minimum width=1.3cm,minimum height=.65cm] (O3) at (-5,0){炉};
  \node[draw,minimum width=1.3cm,minimum height=.65cm] (Z4) at (-2.8,0){SG$_z$};
  \node[draw,minimum width=1.3cm,minimum height=.65cm] (X2) at (.2,.45){SG$_x$};
  \node[draw,minimum width=1.3cm,minimum height=.65cm] (Z5) at (3.2,.75){SG$_z$};
  \draw[->] (O3)--(Z4);
  \draw[->] (Z4.east)--(X2.west) node[midway,above]{$z+$};
  \draw[->] (X2.east)--(Z5.west) node[midway,above]{$x+$};
  \draw[->] (Z5.east)--(5.3,1.05) node[right]{$z+$};
  \draw[->] (Z5.east)--(5.3,.45) node[right]{$z-$};
  \node[left] at (-6.1,0){(c)};
 \end{scope}
\end{tikzpicture}
\bookcaption{序列 Stern--Gerlach 实验：选择一个方向的本征态会破坏另一方向的确定性}
\label{advA6:fig:sequential-sg}
\end{minipage}\end{center}

图 \ref{advA6:fig:sequential-sg}(a) 表明重复测量同一分量结果确定；
(b) 表明 $|z+\rangle$ 对 $S_x$ 是等概率叠加；
(c) 表明选择 $x+$ 后再次测量 $S_z$ 又分成两束。
因此态振幅必须是复数：
\[
 |\vect n,+\rangle=
 \cos\frac\theta2|z+\rangle
 +\e^{\ii\phi}\sin\frac\theta2|z-\rangle,
\]
概率由模方给出。相位 $\phi$ 不影响单次 $S_z$ 概率，却影响其他方向的干涉结果，
这正是实概率模型无法表达的结构。
